\subsection{Space obtained by contraction of a subspace}

Let X be a topological space and let A be a subset of X. Consider the equivalence relation R in X which is the finest for which all elements of A are equivalent: two elements of X are equivalent under this relation if they are equal or if they both belong to A.

DEFINITION 10. --- The quotient space of X by R is denoted X/A and is called the space obtained from X by contracting the subset A.

Denote by \(\rho\colon X\to X/A\) the canonical surjection. Let Y be a topological space and let \(f\colon X\to Y\) be a continuous map. For there to exist a continuous map \(g\colon X/A\to Y\) such that \(g\circ\rho=f\) , it is necessary and sufficient that \(f\) be constant on A.

If the set A is closed (resp. open) in X, the map \(\rho\) is closed (resp. open) and induces a homeomorphism of \(X\setminus A\) onto its image. If A is a closed and quasi-compact subset of X, the map \(\rho\) is proper (TG, I, p. 75, th. 1).

If the set A is empty, \(\rho\) is a homeomorphism. If it is not empty, A is a point of X/A which is called the base point of X/A and is denoted \(s_{\mathrm{X / A}}\) . The space X/A is then identified with the space obtained by attaching to the space \(\{s_{\mathrm{X / A}}\}\) the space X by means of the unique map from A to \(\{s_{\mathrm{X / A}}\}\) .

PROPOSITION 10. --- Let X be a topological space and let A be a subset of X. Suppose that there exists a homotopy \(\sigma: \mathrm{X} \times \mathbf{I} \to \mathrm{X}\) with source \(\mathrm{Id}_{\mathrm{X}}\), constant on \(\mathrm{A} \times \{1\}\) and such that \(\sigma(\mathrm{A} \times \mathbf{I}) \subset \mathrm{A}\). Then the canonical map \(\rho\) from X to X/A is a homotopy equivalence.

Denote by \(f\) the term of \(\sigma\). This is a continuous map from X to X, homotopic to \(\mathrm{Id}_{\mathrm{X}}\). It is constant on A; hence there exists a unique continuous map \(g \colon \mathrm{X} / \mathrm{A} \to \mathrm{X}\) such that \(g \circ \rho = f\). On the other hand, since \(\sigma(\mathrm{A} \times \mathbf{I}) \subset \mathrm{A}\), there exists a map \(\sigma'\) from \(\mathrm{X} / \mathrm{A} \times \mathbf{I}\) into \(\mathrm{X} / \mathrm{A}\) such that \(\sigma'(\rho(x), t) = \rho(\sigma(x, t))\) for all \(x \in \mathrm{X}\) and every \(t \in \mathbf{I}\). By I, p. 19, prop. 10, the map \(\sigma'\) is continuous. It is a homotopy connecting the identity map of \(\mathrm{X} / \mathrm{A}\) to the map \(\rho \circ g\). Consequently, the maps \(g\) and \(\rho\) are homotopy equivalences inverse to each other up to homotopy.

Remarks. --- Let X be a topological space and let A be a subset of X.

\begin{enumerate}
\def\labelenumi{\arabic{enumi})}
\item
  If the canonical map from X onto X/A is a homotopy equivalence, then there exists a homotopy \(\sigma\colon\mathrm{X}\times\mathbf{I}\to\mathrm{X}\) with source \(\mathrm{Id}_{\mathrm{X}}\) which is constant on \(\mathrm{A}\times\{1\}\). But it may happen that none exists satisfying in addition the condition \(\sigma(\mathrm{A}\times\mathbf{I})\subset\mathrm{A}\) (III, p. 322, exerc. 4).
\item
  There may exist a homotopy connecting the identity map of X to a map from X into X that is constant on A without the spaces X and X/A being homotopy equivalent (III, p. 325, exerc. 14).
\item
  Suppose that A is contractible. If the pair (X, A) has the homotopy extension property, then the canonical map from X to X/A is a homotopy equivalence. Indeed, let \(\sigma\colon A \times I \to A\) be a homotopy whose initial map is the identity map of A and whose terminal map is a constant map. There then exists a homotopy \(\sigma'\) with origin \(\mathrm{Id}_{X}\) which extends \(\sigma\). The assertion thus follows from Proposition 10.
\item
  Let X and Y be topological spaces, let A be a nonempty subspace of X, and let B be a subspace of Y. Let \(f: X \to Y\) be a continuous map such that \(f(A) \subset B\). Denote by \(\varphi: X/A \to Y/B\) the continuous map induced by f on the quotients. If f defines a homotopy equivalence of the pair (X, A) onto the pair (Y, B), then the map \(\varphi\) is a homotopy equivalence of the pair (X/A, \(s_{X/A}\)) onto the pair (Y/B, \(s_{Y/B}\)). Let P be a topological space reduced to a single point, let \(\alpha\) and \(\beta\) be the constant maps from A and B into P. Observe that the map \(\varphi\) is identified with the map from \(P \cup_{\alpha} X\) into \(P \cup_{\beta} Y\) induced by f and the identity map of P. The assertion then follows from Corollary 1 of III, p. 249.
\end{enumerate}

